%!TEX program = xelatex
%=============================================================================
% 基于强化学习的动态 TDMA MAC 协议调度优化（bensz-paper / paper-sci-01）
% 编译: xelatex → biber → xelatex → xelatex
%=============================================================================
\documentclass[12pt,a4paper]{article}

%=============================================================================
% 中文支持：ctex 先于 bensz-paper 加载，提供 xeCJK 与中文字体环境
%=============================================================================
\usepackage[UTF8,fontset=fandol]{ctex}

%=============================================================================
% bensz-paper 公共样式（SCI/期刊论文 paper-sci-01 profile）
%=============================================================================
\usepackage{bensz-paper}
\BenszPaperSetup{
  template = paper-sci-01,
  bibliography-style = gb7714-2015,
  line-numbers = true
}

%=============================================================================
% 中文期刊样式覆盖：bensz-paper 默认无章节号；这里恢复 "1 引言 / 1.1 子节" 风格
%=============================================================================
\setcounter{secnumdepth}{3}   % section/subsection/subsubsection 编号；\paragraph 不编号
\titleformat{\section}
  {\bfseries\fontsize{14pt}{18pt}\selectfont}
  {\thesection}{0.6em}{}
\titleformat{\subsection}
  {\bfseries\fontsize{12pt}{16pt}\selectfont}
  {\thesubsection}{0.5em}{}
\titleformat{\subsubsection}
  {\bfseries\normalsize}
  {\thesubsubsection}{0.5em}{}

%=============================================================================
% bensz-paper 未预加载的扩展包（数学 / TikZ / 算法 / 子图 / siunitx / cleveref）
%=============================================================================
\usepackage{mathtools}
\usepackage{amssymb}
\usepackage{amsthm}
\usepackage{subcaption}
\usepackage{float}
\usepackage{siunitx}
\sisetup{detect-all}

\usepackage{tikz}
\usetikzlibrary{arrows.meta, positioning, shapes.geometric}

\usepackage{algorithm}
\usepackage{algpseudocode}
\floatname{algorithm}{算法}
\renewcommand{\listalgorithmname}{算法索引}

%=============================================================================
% 定理类环境
%=============================================================================
\theoremstyle{plain}
\newtheorem{theorem}{定理}[section]
\newtheorem{lemma}[theorem]{引理}
\newtheorem{proposition}[theorem]{命题}
\newtheorem{corollary}[theorem]{推论}
\theoremstyle{definition}
\newtheorem{definition}{定义}[section]
\newtheorem{example}{例}[section]
\theoremstyle{remark}
\newtheorem{remark}{注记}[section]

%=============================================================================
% 参考文献（bensz-paper 已预载 biblatex；vancouver / gb7714-2015 二选一）
%=============================================================================
\addbibresource{references/refs.bib}
\ExecuteBibliographyOptions{url=true, doi=true}

%=============================================================================
% cleveref（必须在 hyperref 之后；bensz-paper 已先载入 hyperref）
%=============================================================================
\usepackage{cleveref}
\crefformat{equation}{式~(#2#1#3)}
\Crefformat{equation}{式~(#2#1#3)}
\crefformat{figure}{图~#2#1#3}
\Crefformat{figure}{图~#2#1#3}
\crefformat{table}{表~#2#1#3}
\Crefformat{table}{表~#2#1#3}
\crefformat{section}{第~#2#1#3~节}
\Crefformat{section}{第~#2#1#3~节}
\crefformat{subsection}{第~#2#1#3~节}
\Crefformat{subsection}{第~#2#1#3~节}
\crefformat{algorithm}{算法~#2#1#3}
\Crefformat{algorithm}{算法~#2#1#3}

%=============================================================================
% 自定义命令
%=============================================================================
\newcommand{\R}{\mathbb{R}}
\newcommand{\N}{\mathbb{N}}
\newcommand{\Z}{\mathbb{Z}}
\DeclareMathOperator*{\argmax}{arg\,max}
\DeclareMathOperator*{\argmin}{arg\,min}
\DeclarePairedDelimiter{\abs}{\lvert}{\rvert}
\DeclarePairedDelimiter{\norm}{\lVert}{\rVert}

%=============================================================================
% 浮动体放置策略：放宽默认阈值，减少大图/大表造成的空白页
%   - topfraction      ↑：允许顶部浮动体占用更多版面
%   - bottomfraction   ↑：允许底部同上
%   - textfraction     ↓：允许页面文字比例更低（少留底）
%   - floatpagefraction↑：浮动页面必须装满到此比例才被接受
%   - topnumber/bottomnumber/totalnumber：放宽每页可放浮动体数量
%=============================================================================
\renewcommand{\topfraction}{0.92}
\renewcommand{\bottomfraction}{0.85}
\renewcommand{\textfraction}{0.06}
\renewcommand{\floatpagefraction}{0.80}
\setcounter{topnumber}{4}
\setcounter{bottomnumber}{4}
\setcounter{totalnumber}{8}

%=============================================================================
\begin{document}

{
\thispagestyle{fancy}
\begin{center}
{\heiti\bfseries\fontsize{15.12}{20}\selectfont 基于强化学习的动态 TDMA MAC 协议调度优化研究\par}
\end{center}
}

\section*{摘要}

\noindent
动态多跳无线网络的时分多址（TDMA）调度需要同时兼顾吞吐量、公平性与拓扑扰动恢复能力，
而传统启发式调度依赖手工权值难以适应拓扑与业务的快速变化，
直接将申请概率替换为深度强化学习（DRL）策略又会破坏已有的安全退避机制并产生饥饿节点。
本文提出\textbf{密度感知服务债务 PPO}（Density-Aware Service-Debt PPO, DASD-PPO）调度框架，
其核心思路是将 PPO 策略以\emph{乘数}形式叠加在启发式申请概率之上，
并在采样阶段引入两跳安全掩膜、服务债务推力、请求预算约束与密度自适应控制，
使学习型策略在拓扑感知归纳偏置下做出公平且鲁棒的逐时隙决策。
DASD-PPO 采用 LSTM Actor--Critic 主干（每节点独立、约 13.8K 参数），
对每个数据时隙单独输出 Bernoulli 申请概率，并配合逐时隙奖励分解解决多维动作的信用归因问题；
仿真器与训练器通过 POSIX 命名管道完成异步在线交互。
在 OMNeT++ 6.3 平台上的对比实验涵盖两个静态主场景（行人 MANET 与 UAV 稀疏拓扑）
与四类动态扰动场景（边切换、桥接断裂、节点重接入、节点永久脱离），
将 DASD-PPO 与 PPO-Base、HC-PPO 及七种启发式/论文启发式近似算法做了同口径对比。
结果显示，DASD-PPO 在 UAV 稀疏拓扑下将饥饿比例从 \(0.5556\) 显著降至 \(\mathbf{0.2500}\)
（全表最低），在 Bridge Break 扰动下取得最小 PDR 衰减，
在 Node Rejoin 恢复阶段取得最高 throughput（\(2.6072\)），
学习型方法整体在 PDR 与 throughput 上数倍优于论文启发式近似基线，
证实在线 RL 调度对动态拓扑具有结构性鲁棒优势。

\vspace{0.5em}
\noindent\textbf{关键词：}时分多址；强化学习；MAC 协议；近端策略优化；空间复用；公平性；服务债务；在线学习


\section{引言}
\label{sec:intro}

无线自组织网络在物联网、车载通信与无人机集群等场景中已成为关键的承载技术，
其 MAC 层的调度策略直接决定网络的吞吐量、时延与节点公平性~\cite{author2024survey}。
时分多址（TDMA）通过将信道划分为固定帧和时隙、为节点分配无冲突的接入机会，
是上述场景中实现确定性低时延的主流方案，并可通过空间复用进一步提高频谱效率~\cite{spatial_reuse}。
然而，节点移动、链路开断与业务异构使可调度时隙的分布持续漂移，
原本针对静态拓扑设计的启发式调度难以同时维持高吞吐与可控的长尾时延。

传统 TDMA 调度算法可分为两类：
\emph{所有者时隙}类方案（plain TDMA 及其增强版本）按固定划分保证无冲突，
但稀疏拓扑下大量时隙闲置而密集区域产生饥饿；
\emph{两跳仲裁与混合分配}类方案（如 Z-MAC~\cite{zmac}、TRAMA~\cite{trama}、STDMA~\cite{stdma}）
通过混合分配/竞争或确定性两跳仲裁缓解上述问题，
但其权值更新仍依赖本地参数手工调优，
在到达过程剧烈波动与拓扑突变下难以兼顾吞吐、公平性与扰动恢复速度。
近年来兴起的基于深度强化学习的 MAC 调度~\cite{drlmac2019,ppoma2023,mnih2015dqn}
为动态调度提供了数据驱动的替代路径。
然而，本文在沿用这一思路构造朴素基线时观察到：
若直接由 PPO 端到端输出申请概率（本文将该朴素基线记为 \textbf{PPO-Base}），
会带来两类新风险：
其一，学习型策略容易破坏仿真器中已有启发式调度器的安全退避先验，
在训练初期反而恶化吞吐与公平性；
其二，对多维 Bernoulli 动作仅给出聚合奖励，无法精确归因每个时隙决策的好坏，
导致策略熵塌缩或学习停滞。
本文正是从修复 PPO-Base 的上述两点不足出发，
逐步引出后述的 DASD-PPO 设计。

本文针对上述问题，提出\textbf{密度感知服务债务 PPO}（DASD-PPO）调度框架。
DASD-PPO 在策略空间中显式引入与拓扑感知一致的归纳偏置，
其核心是\emph{将 PPO 学习信号约束在启发式可行域内，并对长尾未服务节点提供确定性补偿}。
具体而言，DASD-PPO 由四个相互配合的机制组成：
（i）\textbf{乘数模式策略合成与启发式偏离锚定}，
让 Actor 输出乘数因子 \(\alpha\in(0,1)^M\)，
通过 \(p_\text{req}=\text{clamp}(p_\text{heur}\cdot 2\alpha,0,1)\) 合成申请概率，
并以与启发式概率的 \(\ell_2\) 距离作为软惩罚锚定策略；
（ii）\textbf{两跳安全掩膜}，
在采样阶段显式屏蔽上一帧一/二跳邻居占用的时隙，
保留拓扑感知的可行域；
（iii）\textbf{服务债务与请求预算耦合}，
对超过门限连续未发送且队列非空的节点施加确定性动作推力，
并以预算上限防止过申请；
（iv）\textbf{密度感知自适应控制}，
按归一化网络密度 \(\rho_t\) 在掩膜启用/关闭之间平滑插值，
使框架能在稀疏与密集拓扑间统一适配。

为进一步定位上述机制的边际贡献，本文构造了从弱到强的三层学习型算法演进：
\textbf{PPO-Base}（基础逐时隙 Bernoulli 采样）、
\textbf{HC-PPO}（启发式约束 PPO，在 PPO-Base 之上加入启发式偏离软惩罚）、
\textbf{DASD-PPO}（在 HC-PPO 之上再叠加两跳掩膜、服务债务、请求预算与密度自适应）。
三者共用相同的 PPO 训练超参数与训练预算，
其差异完全由前述四个机制贡献。
此外，本文实现了 OMNeT++ 仿真器与外部 PPO 训练器的异步在线训练管道：
仿真器与训练器通过两个 POSIX 命名管道完成状态/动作交换，
任何一端未就绪时另一端自动降级或重连，
为在动态拓扑上做长时序在线学习提供工程基础。

本文的主要贡献可概括为：
\begin{itemize}
    \item \textbf{策略框架}：提出 DASD-PPO 调度算法，
          以乘数模式 + 两跳掩膜 + 服务债务 + 密度自适应四项机制为骨架，
          显式约束 RL 策略不偏离拓扑感知可行域，
          并对长尾节点提供确定性补偿。
    \item \textbf{算法机理}：基于 LSTM Actor--Critic（每节点独立、约 13.8K 参数）
          实现轻量级在线策略；针对多维 Bernoulli 动作引入\emph{逐时隙}信用分配，
          使 Critic 输出 \(M\) 维逐时隙价值，PPO 采用逐时隙比率 \(\times\) 优势函数更新，
          显著提升熵收敛速度与样本效率。
    \item \textbf{系统实现}：构建 OMNeT++ 6.3 仿真器与外部 PPO 训练器的
          POSIX 命名管道双向异步交互管道，支持在线训练与启发式降级，
          能在两种 \(N=12\) 主场景与四种动态扰动场景上稳定运行。
    \item \textbf{完备评估}：在两个静态主场景与四个动态扰动场景上做了 3 个随机种子下的
          十种算法横向对比；结果表明 DASD-PPO 在 UAV 稀疏拓扑下取得\emph{全表最低}的饥饿比例
          （\(0.2500\) vs. HC-PPO 的 \(0.5556\)），
          在 Bridge Break 与 Node Rejoin 扰动下分别取得最小 PDR 衰减与最高 throughput 恢复，
          学习型方法整体相对论文启发式近似（Z-MAC*、TRAMA*、Bandit-Index*）
          PDR 与 throughput 提升数倍。
\end{itemize}

本文其余部分组织如下：
\Cref{sec:related} 回顾相关工作；
\Cref{sec:model} 给出系统模型与优化目标的形式化描述；
\Cref{sec:method} 详细介绍 DASD-PPO 的四个机制与训练流程；
\Cref{sec:exp} 报告静态主场景、动态扰动场景及算法演进消融的实验结果；
\Cref{sec:discussion} 讨论设计权衡与局限；
\Cref{sec:conclusion} 总结全文并展望未来工作。

\section{相关工作}
\label{sec:related}

本节按三条技术线索回顾与本文最直接相关的研究，
并在每条线索末尾点明 DASD-PPO 与现有方法的关键差异。

\subsection{TDMA 调度与空间复用}

固定 TDMA 通过预分配所有者时隙保证无冲突接入，
但在动态拓扑下大量所有者时隙闲置~\cite{author2024survey}。
为提升频谱效率，
图染色与几何冲突图类方法（如 STDMA~\cite{stdma}）
将时隙复用建模为染色问题，
通过两跳安全约束允许相距足够远的节点同时占用同一时隙。
\textbf{Z-MAC}~\cite{zmac} 采用 CSMA/TDMA 混合机制：
非所有者节点在所有者未占用时通过竞争争抢时隙，
实现负载自适应；
\textbf{TRAMA}~\cite{trama} 通过基于优先级的确定性两跳仲裁完成无冲突调度，
并为节能引入信令调度树。
\emph{空间复用与吞吐量关系}的理论分析进一步给出了密集拓扑下可达吞吐的上界~\cite{spatial_reuse}。
然而，上述方法的调度规则结构与关键参数均在部署前离线确定、
运行期不随业务到达分布与拓扑变化做在线自适应，
因此在拓扑突变与到达过程剧烈波动场景下仍难以兼顾吞吐与长尾时延。
本文不试图替换其中任意一种协议状态机，
而是在 PPO 学习框架内显式引入\emph{两跳安全掩膜}与\emph{服务债务}两项归纳偏置，
将传统调度的拓扑感知先验作为可学习策略的边界约束。

\subsection{强化学习驱动的 MAC 调度}

用学习方法求解 MAC 决策已积累了若干代表性工作。
其中 DLMA~\cite{drlmac2019} 与 PPOMA~\cite{ppoma2023} 面向单节点共享一条异构信道的二元接入决策，
分别以 DQN 与 PPO 为策略骨干；
Dutta 等~\cite{bandit} 则将每节点的时隙选择建模为分布式多臂老虎机，
作为非深度学习的轻量替代
（MAB 在更广义无线决策中的综述可见~\cite{mab_5g}）。
本文沿用其中 PPO~\cite{schulman2017ppo}、GAE~\cite{schulman2016gae} 与 LSTM~\cite{lSTM_seq} 等成熟组件，
但把目标场景从"单节点二元接入"扩展为"多节点 \(\times\,M\) 维时隙申请 + 空间复用 + 公平性"
的分布式 TDMA 调度。

然而，DLMA/PPOMA 类方法面向的均是\emph{单节点 \(\times\) 单信道}的接入博弈，
其策略空间仅含"发/不发"二元动作，
与本文关注的\emph{分布式多节点 \(\times M\) 维时隙申请 + 空间复用 + 公平性}调度问题在尺度与结构上均不重合。
据我们所知，
将 PPO 策略与完整三阶段 TDMA 调度框架、空间复用约束与服务债务机制\emph{统一在同一在线学习管线}下的代表性公开实现仍属空白。
直接把上述二元决策方法移植到多时隙申请场景时，
还会暴露两类共性缺陷：
\emph{其一}，由策略直接输出申请概率会覆盖已有启发式调度器的退避机制与安全先验，
易在训练初期破坏可行域、引发过申请；
\emph{其二}，\(M\) 维 Bernoulli 申请动作只对应一个聚合奖励，
单时隙决策无法获得精确信用归因，
PPO 的策略熵在小批量更新下容易塌缩。
DASD-PPO 通过\emph{乘数模式策略合成}（保留启发式作为先验基准）与\emph{逐时隙信用分配}
（Critic 输出 \(M\) 维价值、PPO 损失逐时隙独立计算）同时回应上述两个缺陷。

\subsection{公平性与队列稳定性}

公平性是多用户共享信道场景中与吞吐量同等重要的指标。
Jain 等~\cite{jain1984} 定义的 Jain 公平指数已成为评价资源分配公平性的标准工具，
本文沿用其在最后 100 帧的局部值与全过程趋势两种统计。
\textbf{Lyapunov drift 方法}~\cite{neely2010stochastic} 提供了在保证队列稳定的同时最大化效用的形式化框架，
但其拉格朗日乘子需要在线估计且对环境噪声敏感。
本文在工程实验中也曾尝试以拉格朗日变体显式约束饥饿，
然其约束信号在仿真中容易饱和到 \(1.0\)（见 \Cref{sec:discussion}）；
DASD-PPO 选择以\emph{服务债务}的二值指示量驱动确定性动作推力，
辅以请求预算上限抑制过申请，
在工程实现上更稳定，
同时保留了 Lyapunov 思想中"将公平性作为可调节项叠加到效用最大化"的结构。

\section{系统模型与问题描述}
\label{sec:model}

\subsection{网络模型}

考虑一个由 \(N\) 个节点构成的多跳无线自组织网络，
拓扑用无向图 \(\mathcal{G}_t = (\mathcal{V},\mathcal{E}_t)\) 表示，
其中 \(\mathcal{V}=\{0,1,\ldots,N-1\}\)，
\(\mathcal{E}_t\) 为帧 \(t\) 时刻的活跃链路集合（允许随时间变化）。
节点 \(i\) 的一跳邻居集合记为 \(\mathcal{N}_i^{(1)}\)，
两跳邻居集合记为 \(\mathcal{N}_i^{(2)}=\bigcup_{j\in\mathcal{N}_i^{(1)}}\mathcal{N}_j^{(1)}\setminus\{i\}\)。
每个节点维护一个分组发送队列 \(Q_{i,t}\)，
新到达分组按其优先级（低/高/关键）入队。

\subsection{三阶段 TDMA 帧结构}

每个 TDMA 帧包含 \(M\) 个数据时隙，并按以下三个阶段顺序执行：
\begin{enumerate}
    \item \textbf{请求阶段（Request Phase, RTS）}：
        每节点在其专属微时隙广播 \texttt{TDMAGrantRequest}，
        携带对每个数据时隙的申请概率 \(p_{i,m,t}\in[0,1]\)。
    \item \textbf{回复阶段（Reply Phase, CTS）}：
        节点根据收到的 RTS 做出授权决策（+1 同意 / -2 拒绝），
        广播 \texttt{TDMAGrantReply}；
        所有 CTS 汇总后构建本帧的占用表 \(\mathrm{occ}_t[m]\)。
    \item \textbf{数据阶段（Data Phase）}：
        持有 \(\mathrm{mySlots}_{i,t}[m]=1\) 的节点在时隙 \(m\) 发送其优先级最高的队首分组，
        阶段结束时向 RL 训练器推送状态特征并读取下一帧动作。
\end{enumerate}
策略 \(\pi_\theta(\cdot\mid s_t)\) 输出的是每帧每节点的申请概率向量
\(p_{\mathrm{req},i,t}\in[0,1]^M\)，
最终送入仿真的为对 \(p_{\mathrm{req},i,t}\) 独立 Bernoulli 采样得到的二值动作
\(a_{i,t}\in\{0,1\}^M\)。

\subsection{空间复用约束}

时隙 \(m\) 可被多个节点同时占用，
当且仅当任意两个占用节点 \(i,j\) 满足两跳安全条件 \(\mathrm{dist}_t(i,j)\ge 3\)。
本文将该条件作为 DASD-PPO 中两跳掩膜的依据
（\Cref{subsec:mask}）。

\subsection{优化目标}
\label{subsec:objective}

记节点 \(i\) 在帧 \(t\) 时隙 \(m\) 的逐时隙结果指示
\(y_{i,m,t}\in\{0,1,2\}\)：
\(0\) 表示未申请，\(1\) 表示发送成功，\(2\) 表示发生碰撞或拒绝。
定义节点 \(i\) 在帧 \(t\) 的瞬时吞吐与公平性指标：
\begin{equation}
\mathrm{TX}_{i,t} = \sum_{m=1}^{M}\mathbf{1}[y_{i,m,t}=1],
\qquad
J_t = \frac{\left(\sum_i \mathrm{TX}_{i,t}\right)^2}{N\cdot\sum_i \mathrm{TX}_{i,t}^2},
\label{eq:tx-jain}
\end{equation}
其中 \(J_t\in[1/N,1]\) 为 Jain 公平指数~\cite{jain1984}。
长期优化目标可形式化为
\begin{equation}
\max_{\pi_\theta}\;
\mathbb{E}_{\pi_\theta}\!\left[
\sum_{t=0}^{T}\gamma^{t}
\left( \sum_i \mathrm{TX}_{i,t} - \beta\cdot \sum_i \mathrm{Coll}_{i,t} \right)
\right]
\;\;\text{s.t.}\;
\liminf_{t\to\infty} J_t \ge J_{\min},
\;\;
\bar{Q}_{i,t} \text{ 有界}.
\label{eq:problem}
\end{equation}
其中 \(\mathrm{Coll}_{i,t}\) 为本帧碰撞次数，
\(\gamma\in(0,1)\) 为折扣因子，
\(\beta\) 为碰撞惩罚权重，
约束保证长期公平性与队列稳定性。
\eqref{eq:problem} 在动态拓扑下没有闭式解，
本文以 PPO 的逐时隙信用分解将其松弛为可在线优化的代理目标
（\Cref{subsec:credit}）。

\section{方法：DASD-PPO 调度算法}
\label{sec:method}

\subsection{总览与算法演进}
\label{subsec:overview}

DASD-PPO 是一个以 PPO 为骨干、
叠加四项拓扑感知归纳偏置的在线 TDMA 调度算法。
其设计动机源于一对观察：
\emph{一方面}，传统 MAC 调度文献长期积累的两跳安全约束、所有者优先与
失败时隙退避等规则~\cite{stdma,zmac,trama}
是经过数十年验证的高质量安全调度先验，
直接抛弃这些先验、要求 RL 从零样本中重新学到一致的安全行为，
既样本效率低、又在训练初期容易破坏可行域；
\emph{另一方面}，固定权值的启发式调度无法对长尾节点的长时间未服务做出显式补偿，
也无法跨拓扑密度自适应调整空间复用的激进程度。
因此，DASD-PPO 将 RL 策略\emph{约束在传统启发式可行域之内}，
并\emph{显式}为长尾节点提供补偿信号，
同时按网络密度动态调整安全约束的强度。

为便于讨论上述四项机制的边际贡献，
本文在内部构造了从弱到强的四层算法演进
（见 \Cref{fig:algo-evolution} 与 \Cref{tab:algo-naming}）：
\textbf{PPO-Base} 仅保留逐时隙 Bernoulli 采样的基础策略空间；
\textbf{HC-PPO} 在 PPO-Base 之上加入启发式偏离软惩罚，
将策略锚定在启发式邻域；
\textbf{DASD-PPO w/o Mask} 在 HC-PPO 之上进一步引入服务债务、请求预算与密度自适应控制三项机制，
但\emph{不}启用两跳安全掩膜——它是为隔离掩膜边际贡献而引入的消融对照变体；
\textbf{DASD-PPO} 在 DASD-PPO w/o Mask 之上再启用两跳安全掩膜与配套的 \(\alpha=0.75\) 安全槽偏置，
是本章设计意图最完整的主方法。
四者共享相同的状态/动作空间与 PPO 训练超参数，
在 \Cref{sec:exp} 中构成四级模块化消融。

\begin{table}[H]
\centering
\caption{本文学习型算法命名与实验代号对照}
\label{tab:algo-naming}
\begin{tabular}{@{}lll@{}}
\toprule
\textbf{论文名}（缩写） & \textbf{中文全称} & \textbf{实验代号 / CSV 列名} \\
\midrule
PPO-Base          & 基础 PPO 调度器                        & \texttt{ppo\_baseline} \\
HC-PPO            & 启发式约束 PPO                         & \texttt{B} \\
DASD-PPO w/o Mask & 密度感知服务债务 PPO（无掩膜消融变体） & \texttt{B\_debt\_adaptive} \\
DASD-PPO          & 密度感知服务债务 PPO                   & \texttt{B\_masked\_debt\_adaptive} \\
\bottomrule
\end{tabular}
\end{table}

\begin{figure}[H]
\centering
\begin{tikzpicture}[
    node distance=8mm and 10mm,
    box/.style={
        draw, rounded corners=2.4pt, align=center,
        inner sep=3.5pt, minimum width=30mm, minimum height=12mm,
        font=\small
    },
    main/.style={box, fill=red!8, very thick},
    sub/.style={box, fill=blue!4},
    abl/.style={box, fill=orange!12, dashed},
    arr/.style={-{Latex[length=2.2mm]}, thick},
    note/.style={font=\scriptsize, align=center}
]
\node[sub] (m0) {PPO-Base \\ \scriptsize 逐时隙 Bernoulli 采样};
\node[sub, right=of m0] (m1) {HC-PPO \\ \scriptsize + 启发式偏离惩罚};
\node[abl, right=of m1] (m2) {DASD-PPO w/o Mask \\ \scriptsize + 服务债务 + 请求预算 \\ \scriptsize + 密度自适应};
\node[main, right=of m2] (m3) {DASD-PPO \\ \scriptsize + 两跳安全掩膜};

\draw[arr] (m0) -- node[above, note] {KL 形锚定} (m1);
\draw[arr] (m1) -- node[above, note] {公平性 + 自适应} (m2);
\draw[arr] (m2) -- node[above, note] {拓扑感知掩膜} (m3);
\end{tikzpicture}
\caption{算法演进示意图：
PPO-Base \(\rightarrow\) HC-PPO \(\rightarrow\) DASD-PPO w/o Mask \(\rightarrow\) DASD-PPO。
箭头标注每一阶段相对前代引入的关键机制。}
\label{fig:algo-evolution}
\end{figure}

后续 \Cref{subsec:srl} 至 \Cref{subsec:training} 按上述顺序依次给出 DASD-PPO 的各个组成部分，
最后在 \Cref{alg:dasd-ppo} 中给出整体的一帧训练流程。

\subsection{状态、动作与奖励}
\label{subsec:srl}

\paragraph{状态空间。}
每节点每帧观测的状态向量维度为 \(5M+10+N\)（默认 \(M=10,N=12\) 时共 \(72\) 维），
由四类特征拼接而成：
\begin{itemize}
    \item \textbf{时隙感知（\(3M\) 维）}：
          Bown（自身上一帧占用状态，\(M\) 维）；
          T2hop（两跳占用标识 \(+\) 最小跳数归一化，\(2M\) 维）。
    \item \textbf{队列与流量（10 维标量）}：
          队列长度 \(Q_{i,t}\)、到达率 EWMA、服务债务计数、邻居利用率等。
    \item \textbf{公平性与节点 ID（\(N\) 维）}：节点 ID 的 one-hot 编码。
    \item \textbf{策略上下文（\(2M\) 维）}：
          上一帧申请概率向量 \(P_{t-1}\)（\(M\) 维）、
          本帧仿真器侧启发式基准概率 \(p_{\mathrm{heur},t}\)（\(M\) 维）。
\end{itemize}
\(p_{\mathrm{heur},t}\) 作为乘数模式下的基准参考，让 RL 能够精确学习相对启发式的偏移。

\paragraph{动作空间。}
Actor 输出乘数因子 \(\alpha_t\in(0,1)^M\)，
申请概率按乘数模式合成：
\begin{equation}
p_{\mathrm{req},t} \;=\; \mathrm{clamp}\!\left( p_{\mathrm{heur},t}\cdot 2\alpha_t,\;0,\;1\right).
\label{eq:multiplier}
\end{equation}
\(\alpha_t\approx 0.5\) 时等效纯启发式；
Actor 输出层零初始化保证 Sigmoid 输出在训练起点为 \(0.5\)，
即学习从启发式策略平稳启动，
避免训练初期对稳定时隙做出激进偏离。
最终送入仿真的为对 \(p_{\mathrm{req},t}\) 做独立 Bernoulli 采样得到的二值动作 \(a_t\in\{0,1\}^M\)。

\paragraph{奖励信号。}
为支撑多维动作的精确信用分配，
仿真器在每帧逐时隙向训练器上报每个时隙的发送结果
（未申请 / 发送成功 / 发生碰撞，三选一），
训练器据此计算逐时隙奖励：
\begin{equation}
r_{m,t} = \begin{cases}
        +1, & y_{m,t}=1\;\text{（发送成功）}, \\
        -1, & y_{m,t}=2\;\text{（发生碰撞）}, \\
        \phantom{+}0, & y_{m,t}=0\;\text{（未申请）}.
\end{cases}
\label{eq:slot-reward}
\end{equation}
此外，为对齐 \(t\) 时刻动作的实际反馈，
本文引入一个\emph{时序对齐缓冲}：
帧 \(t\) 上报的 \(r_{m,t}\) 与帧 \(t-1\) 的动作 \(a_{m,t-1}\) 配对后再写入回放缓冲，
修复因果时序错位。
最后对每节点维护 EWMA 奖励基线（衰减 \(0.95\)），
存入 buffer 的奖励为差分奖励 \(r^{\mathrm{shaped}}_{m,t}=r_{m,t}-\bar r_{m,t}\)，
减小方差、突出改进信号。

\subsection{LSTM Actor--Critic 网络}
\label{subsec:net}

DASD-PPO 采用轻量级 LSTM Actor--Critic 架构，
每节点独立维护一组参数与优化器（非参数共享）；
单节点网络总参数量约 \(13.8\)K。
之所以选用如此轻量的网络规模，
是因为本文每帧只能采集到 \(N\) 个节点 \(\times M\) 个时隙的局部观测样本，
样本规模有限；
若网络过参数化，PPO 在小批量更新下梯度信号易被噪声淹没、
策略熵塌缩，
难以稳定收敛。
前向过程为：
\begin{equation}
    \begin{aligned}
        h_t      &= \mathrm{LSTM}(x_t,\;h_{t-1};\;\theta_{\mathrm{enc}}),\quad h_t\in\R^{32},\\
        \alpha_t &= \sigma(W_a h_t + b_a),\quad \alpha_t\in(0,1)^M,\\
        V_t      &= W_c h_t + b_c,\quad V_t\in\R^M,
    \end{aligned}
    \label{eq:network}
\end{equation}
其中 \(\sigma(\cdot)\) 为 Sigmoid 函数。
\(W_a,b_a\) 零初始化保证学习起点等效纯启发式（\(\alpha_t\equiv 0.5\)）；
Critic 头输出 \(M\) 个逐时隙价值估计 \(V_{m,t}\)，
为后述逐时隙信用分配提供基线函数。

\subsection{逐时隙信用分配}
\label{subsec:credit}

传统聚合奖励 \(r=\sum_m r_{m,t}\) 与单标量价值估计配合时，
\(M\) 个独立 Bernoulli 动作只能共享同一个 advantage 信号，
PPO 无法区分哪个时隙决策好/坏。
DASD-PPO 改为对每个时隙独立计算优势：
\begin{equation}
\hat A_{m,t} = \sum_{l=0}^{L-1}(\gamma\lambda)^{l}\,\delta_{m,t+l},
\qquad
\delta_{m,t} = r^{\mathrm{shaped}}_{m,t} + \gamma V_{m,t+1} - V_{m,t},
\label{eq:gae-slot}
\end{equation}
PPO 策略损失改为逐时隙形式：
\begin{equation}
\mathcal{L}_{\mathrm{policy}}(\theta) =
- \frac{1}{M}\sum_{m=1}^{M}
\min\!\left(
\rho_{m,t}(\theta)\,\hat A_{m,t},\;
\mathrm{clip}(\rho_{m,t}(\theta),1-\epsilon,1+\epsilon)\,\hat A_{m,t}
\right),
\label{eq:ppo-slot}
\end{equation}
其中比率 \(\rho_{m,t}(\theta)=\pi_\theta(a_{m,t}\mid s_t)/\pi_{\theta_{\mathrm{old}}}(a_{m,t}\mid s_t)\)。
实验显示，该改造使策略 entropy 收敛速度提升约 \(105\times\)，
L\_actor 从死寂（\(\approx 0\)）变为活跃（\(+0.0076\)），
证实逐时隙信用分配有效解决了多动作维度的归因问题（详见 \Cref{subsec:limitations}）。

\subsection{启发式偏离锚定（HC-PPO 模块）}
\label{subsec:hc}

PPO-Base 的纯学习策略容易在训练初期对稳定时隙做出激进偏离，
破坏仿真器侧已有启发式的安全调度先验。
HC-PPO 模块在 PPO 损失中加入与启发式概率向量的二阶距离软惩罚：
\begin{equation}
\mathcal{L}_{\mathrm{heur}}(\theta) \;=\; \lambda_{\mathrm{heur}}\cdot
   \mathbb{E}_{t}\!\left[\,\norm{p_{\mathrm{req},t}-p_{\mathrm{heur},t}}_2^{2}\,\right],
\qquad \lambda_{\mathrm{heur}}\in[0.01,0.02].
\label{eq:heur-deviation}
\end{equation}
\eqref{eq:heur-deviation} 等价于在策略空间施加以启发式为先验的 KL 形锚定，
使学习型策略在保持探索能力的同时不脱离安全调度的可行域；
HC-PPO 因此作为 DASD-PPO 的内层基底，
后续两跳掩膜、服务债务等机制均叠加在 HC-PPO 之上。

\subsection{两跳安全掩膜}
\label{subsec:mask}

仅靠软惩罚不足以避免学习型策略采样到与近邻冲突的时隙。
DASD-PPO 在采样阶段引入一个\emph{硬约束算子}，
\emph{显式屏蔽}上一帧一跳/二跳邻居已占用的时隙：
\begin{equation}
\tilde p_{\mathrm{req},m,t} \;=\;
\begin{cases}
p_{\mathrm{req},m,t}, & \text{若 } m\in\mathcal{S}_{\mathrm{safe},t-1},\\
0,                    & \text{否则},
\end{cases}
\label{eq:mask}
\end{equation}
同时安全时隙的 Actor 输出初始化为 \(\alpha_m=0.75\)
（等效初始申请概率为 \(1.5\,p_{\mathrm{heur},m}\)），
鼓励学习期对安全候选时隙做更积极的尝试。
两跳掩膜将策略可行域显式收缩到与近邻无冲突的子空间，
保留了 \Cref{sec:related} 中传统两跳仲裁方法的核心拓扑感知机制，
但其激活与否由后述密度自适应控制按当前 \(\rho_t\) 决定。

\subsection{服务债务与请求预算}
\label{subsec:debt}

两跳掩膜虽能避免冲突，
却无法补偿"明明有队列但长期未发送"的长尾节点。
为此 DASD-PPO 引入\emph{服务债务}指示量：
\begin{equation}
D_{i,t} \;=\; \mathbf{1}\!\left[\,f_{i,t}>\tau_D\;\land\;Q_{i,t}>0\,\right],
\qquad \tau_D=5,
\label{eq:debt-indicator}
\end{equation}
其中 \(f_{i,t}\) 为节点 \(i\) 自上次成功发送以来的连续未发送帧数。
当 \(D_{i,t}=1\) 时，在安全时隙上施加确定性动作推力：
\begin{equation}
\tilde\alpha_{i,m} \;=\; \mathrm{clamp}\!\left(\alpha_{i,m}+\beta_D,\;0,\;1\right),
\qquad \beta_D=0.15,
\label{eq:debt-boost}
\end{equation}
并在成功发送时给予奖励加成 \(r_{\mathrm{debt}}=\gamma_D\cdot\mathbf{1}[\text{成功}]\)（\(\gamma_D=0.75\)）。

为防止债务驱动的动作推力造成新的过申请，
设节点 \(i\) 在帧 \(t\) 的有效安全时隙数为 \(\abs{\mathcal{S}_{\mathrm{safe},i,t}}\)，
引入\emph{请求预算}约束：
\begin{equation}
\sum_{m\in\mathcal{S}_{\mathrm{safe},i,t}}\tilde\alpha_{i,m}
\;\le\; B_{\mathrm{req}}\cdot\kappa_B,
\qquad B_{\mathrm{req}}=5.0,\;\kappa_B\in[0.9,1.1].
\label{eq:budget}
\end{equation}
\(\kappa_B\) 由近期申请成功率与队列增长动态缩放；
超过上限时按比例下调全部 \(\tilde\alpha_{i,m}\)，
保证债务推力不会无限放大。

\subsection{密度感知自适应控制}
\label{subsec:density}

两跳掩膜在不同密度下的代价并不对称：
稀疏拓扑下安全集合较大，掩膜能精确收缩可行域而不损失候选；
密集拓扑下安全集合萎缩，过严的掩膜会令策略无候选可选。
为此 DASD-PPO 引入由网络密度驱动的自适应控制机制。
设网络在帧 \(t\) 的活跃节点数为 \(n_{\mathrm{act}}\)、活跃边数为 \(e_{\mathrm{act}}\)，
归一化密度为
\begin{equation}
\rho_t \;=\; \frac{e_{\mathrm{act}}}{n_{\mathrm{act}}(n_{\mathrm{act}}-1)/2}.
\label{eq:density}
\end{equation}
当 \(\rho_t<\rho_{\mathrm{sparse}}=0.30\) 时\emph{启用}两跳掩膜，
并保持 \(\alpha\) 向 \(0.75\) 偏置；
当 \(\rho_t>\rho_{\mathrm{dense}}=0.45\) 时\emph{关闭}掩膜，
并令 \(\alpha\) 平滑回退至 \(0.5\)；
中间区间按 \(\rho_t\) 线性插值。
注意：服务债务侧的动作推力与还债奖励\emph{不}受密度因子衰减，
以保证长尾节点在任意密度下都能得到补偿。

\subsection{在线 PPO 训练流程}
\label{subsec:training}

仿真器与训练器之间通过两条命名管道完成异步双向通信：
\emph{状态管道}由仿真器写入、训练器读取，每帧每节点一条 JSON 记录；
\emph{动作管道}由训练器写入、仿真器读取，每帧一条 JSON 记录。
两条管道均采用非阻塞模式：
若训练器未运行，仿真器自动回退到启发式策略并周期性重试连接；
反之，训练器可独立完成 RL 更新。
完整一帧训练流程如 \Cref{alg:dasd-ppo} 所示。

\begin{algorithm}[!htbp]
\caption{DASD-PPO 一帧训练流程}
\label{alg:dasd-ppo}
\begin{algorithmic}[1]
\Require 状态 \(s_t\)；启发式概率 \(p_{\mathrm{heur},t}\)；安全时隙集 \(\mathcal{S}_{\mathrm{safe},t}\)；
         网络密度 \(\rho_t\)；服务债务向量 \(D_t\)；预算 \(B_{\mathrm{req}}\)；更新间隔 \(K=128\)
\State Actor 前向：\(\alpha_t \gets \mathrm{Actor}_\theta(s_t)\)
\If{\(\rho_t<\rho_{\mathrm{sparse}}\)} \Comment{稀疏拓扑}
    \State 启用两跳掩膜，按式~\eqref{eq:mask}\,（两跳安全掩膜）得 \(\tilde p_{\mathrm{req},t}\)
\ElsIf{\(\rho_t>\rho_{\mathrm{dense}}\)} \Comment{密集拓扑}
    \State 关闭掩膜，\(\alpha_t \gets 0.5 + 0.5(\alpha_t-0.5)\cdot w_{\mathrm{dense}}\)
\Else \Comment{过渡区}
    \State 按 \(\rho_t\) 在掩膜与开放之间线性插值
\EndIf
\State 按式~\eqref{eq:debt-boost}\,（服务债务推力）对债务节点的 \(\tilde\alpha_{i,m}\) 加 \(\beta_D\)
\State 若 \(\sum\tilde\alpha_{i,m}\) 超过式~\eqref{eq:budget}\,（请求预算）规定的上限，则按比例下调
\State 按式~\eqref{eq:multiplier}\,（乘数模式策略合成）得 \(p_{\mathrm{req},t}\)，
       对每个时隙做 Bernoulli 采样得动作 \(a_t\)
\State 将 \(a_t\) 送入仿真器执行，收集逐时隙奖励
       \(r_t\in\{-1,0,+1\}^M\) 与还债奖励 \(r_{\mathrm{debt}}\)
\State 将 \((s_{t-1},a_{t-1},r_t+r_{\mathrm{debt}},s_t)\) 经时序对齐后写入回放缓冲
\If{累计样本达到 \(K\)}
    \State 按式~\eqref{eq:gae-slot}\,（逐时隙 GAE）计算优势 \(\hat A_{m,t}\)
    \State 按式~\eqref{eq:ppo-slot}\,（逐时隙 PPO 损失）与式~\eqref{eq:heur-deviation}\,（启发式偏离锚定）
           做小批量 PPO 更新
    \State 衰减学习率 \(\eta \gets \eta\cdot\gamma_{\mathrm{decay}}\)
\EndIf
\end{algorithmic}
\end{algorithm}

\section{实验评估}
\label{sec:exp}

本节按以下顺序组织：
\Cref{subsec:setup} 给出仿真平台、对比算法、评价指标与超参数；
\Cref{subsec:static-results} 报告两个静态主场景下十一种算法的横向对比；
\Cref{subsec:dynamic-results} 报告四种动态扰动场景下的按帧趋势与三阶段统计；
\Cref{subsec:ablation} 通过 PPO-Base \(\rightarrow\) HC-PPO \(\rightarrow\) DASD-PPO w/o Mask \(\rightarrow\) DASD-PPO 的四级演进对比，
分别定位\emph{启发式锚定}、\emph{服务债务/密度自适应/请求预算}以及\emph{两跳掩膜}三组机制的边际贡献。

\subsection{实验设置}
\label{subsec:setup}

\subsubsection{仿真平台与拓扑}

实验基于 OMNeT++ 6.3 实现~\cite{omnetpp}，
通过 POSIX 命名管道与独立运行的 PPO 训练器进程做异步双向通信。
本文主表采用两个 MANET 主场景：
\texttt{N12\_grid\_manet\_pedestrian}（行人移动模型，\(N=12\) 节点，\(M=10\) 时隙）
与 \texttt{N12\_grid\_manet\_uav\_sparse}（无人机灵感稀疏拓扑，
平均度数显著低于行人场景）。
动态扰动评估在行人主场景的四个变体上进行：
\textbf{Edge Toggle}（链路开断切换）、
\textbf{Bridge Break}（关键桥接链路断开）、
\textbf{Node Rejoin}（已下线节点重新加入）、
\textbf{Node Dropout}（节点永久脱离）。
扰动起始帧固定为 \(\tau_{\mathrm{perturb}}=20000\)，
恢复段起始帧为 \(\tau_{\mathrm{recovery}}=40000\)；
Node Dropout 场景按设计不存在恢复段，
其恢复列在统计上以 N/A 标注。

\subsubsection{对比算法}

本文共比较十一种算法，可分为三组：

\paragraph{学习型算法（内部演进）}
PPO-Base、HC-PPO、\textbf{DASD-PPO w/o Mask}（DASD-PPO 关闭两跳掩膜的消融变体）、DASD-PPO，
四者的差异已在 \Cref{tab:algo-naming} 与 \Cref{fig:algo-evolution} 给出。
所有学习型算法均采用同一组 PPO 训练超参数与训练预算。
其中 DASD-PPO w/o Mask 与 DASD-PPO 的唯一差异是 \Cref{subsec:mask} 描述的两跳安全掩膜与配套的 \(\alpha=0.75\) 安全槽偏置，
其它机制（启发式偏离锚定、服务债务推力、请求预算、密度自适应控制）完全相同；
其作用是为 \Cref{subsec:ablation} 提供"掩膜单独贡献"的边际证据。

\paragraph{启发式与 STDMA 强基线}
\textbf{plain TDMA}（仅按所有者时隙分配，无空间复用）；
\textbf{Heuristic}（本仓库实现的贪心启发式申请策略）；
\textbf{Greedy-STDMA-2hop}（两跳安全的贪心空间复用 TDMA~\cite{stdma}）；
\textbf{Traffic-Adaptive TDMA}（所有者时隙保留加基于队列长度的机会复用）。

\paragraph{论文启发式近似（paper-inspired）}
\textbf{Z-MAC*}、\textbf{TRAMA*}、\textbf{Bandit-Index*} 分别为对
Z-MAC~\cite{zmac}、TRAMA~\cite{trama} 与多臂老虎机索引调度~\cite{bandit}
的本地近似实现。星号注明此类方法\emph{未}完整复现原协议状态机，
仅按其核心机制（混合分配/竞争、确定性两跳仲裁、UCB 索引）做轻量化实现，
作为论文文献基线的下界参考。

\subsubsection{评价指标}

静态主场景关注：
\textbf{Goodput/slot}~\(\uparrow\)（单时隙有效吞吐）、
\textbf{PDR}~\(\uparrow\)（分组投递率）、
\textbf{MAC 效率}~\(\uparrow\)、
\textbf{Queue slope}~\(\downarrow\)（队列长度斜率，反映长期稳定性）、
\textbf{\(W_t\) P95}~\(\downarrow\)（队首等待时间 95 分位）、
\textbf{Jain 尾 100 帧}~\(\uparrow\)（最后 100 帧的发送公平指数）、
\textbf{Starvation ratio}~\(\downarrow\)（饥饿节点比例）、
\textbf{有效服务率}~\(\uparrow\)~\(\big({=}\,(1-\mathrm{Starv})\times\mathrm{Goodput}\big)\)
（衡量"既有人发又发得多"的复合指标，
用于消除 Jain 在低吞吐场景下退化为"大家一起穷"型伪公平的偏差）。
动态扰动评估对每个场景将仿真切分为
\(\mathrm{pre}/\mathrm{perturb}/\mathrm{recovery}\) 三阶段，
分别报告\textbf{每帧 PDR} 与\textbf{每帧 throughput} 的平均值。
所有数值取 3 个独立随机种子的平均。

\subsubsection{超参数与训练预算}

PPO 学习率 \(\eta=5\times 10^{-4}\)（每次更新后按 \(0.9995\) 指数衰减），
折扣因子 \(\gamma=0.95\)、GAE 参数 \(\lambda=0.95\)、PPO clip 参数 \(\epsilon=0.2\)、
熵系数 \(c_{\mathrm{ent}}=0.05\)（自适应：下限 \(0.05\)、上限 \(0.10\)，
按 entropy 值在 \([0.45,0.55]\) 区间线性插值）。
LSTM 隐层维度 \(32\)（单层）；
每节点独立优化器、独立 \(\sim 13.5\)K 参数。
正式实验采用 3 个随机种子，
仿真时长为 \(20{,}000\) 帧，
所有学习型方法均要求 PPO 更新数达到 \(400\) 次目标值；
静态主场景实验 \(60/60\) 次运行完整收尾、\(36/36\) 次学习型实验达到目标 PPO 更新次数；
动态扰动实验覆盖 \(4\) 场景 \(\times\, 9\) 方法 \(\times\, 3\) seeds 共 \(108\) 次运行
（其中新增的 DASD-PPO w/o Mask 变体补跑 \(12\) 次，与原 \(96\) 次合并）。

\subsection{静态主场景结果}
\label{subsec:static-results}

\Cref{tab:static-pedestrian} 与 \Cref{tab:static-uav-sparse}
分别给出行人与 UAV 稀疏两个主场景下十一种算法的对比结果。
所有数值直接取自正式实验 summary。

\paragraph{关于 Jain 尾 100 的正确解读。}
Jain 指数 \eqref{eq:tx-jain} 在分母含 \(\sum_i\mathrm{TX}_i^2\)：
当所有节点发送量都同步趋近零时 \(J\to 1\)，即\emph{"大家一起穷"也会被算成完美公平}。
表格中 Traffic-Adaptive TDMA（Jain \(0.890\)、Starvation \(0.917\)）
与 Bandit-Index*（Jain \(0.842\)、Starvation \(0.694\)）即落在该退化区。
因此本文不单看 Jain，而以\emph{有效服务率}~\((1-\mathrm{Starv})\times\mathrm{Goodput}\)
作复合排序，
同时奖励"参与节点多"与"参与节点发得多"。

\begin{table}[!htbp]
\centering
\caption{行人主场景（\texttt{N12\_grid\_manet\_pedestrian}）静态对比，3 seeds 平均}
\label{tab:static-pedestrian}
\setlength{\tabcolsep}{4pt}
\resizebox{\textwidth}{!}{%
\begin{tabular}{@{}lrrrrrrrr@{}}
\toprule
\textbf{方法} &
  \textbf{Goodput/slot}\(\uparrow\) &
  \textbf{PDR}\(\uparrow\) &
  \textbf{MAC eff.}\(\uparrow\) &
  \textbf{Queue slope}\(\downarrow\) &
  \textbf{\(W_t\) P95}\(\downarrow\) &
  \textbf{Jain 尾 100}\(\uparrow\) &
  \textbf{Starv.}\(\downarrow\) &
  \textbf{Eff. Svc.}\(\uparrow\) \\
\midrule
plain TDMA            & 0.0628 & 0.0308 & 0.0524 & 18.2623 & 18716.3 & 0.1069 & 0.6944 & 0.0192 \\
Heuristic             & 0.1473 & 0.0722 & 0.0123 & 18.3855 & 18425.5 & 0.5416 & 0.7778 & 0.0327 \\
Greedy-STDMA-2hop     & 0.1608 & 0.0789 & 0.0136 & 17.9578 & 18343.5 & 0.6849 & 0.8333 & 0.0268 \\
Traffic-Adaptive TDMA & 0.0267 & 0.0131 & 0.0022 & 19.9179 & 18757.5 & 0.8900 & 0.9167 & 0.0022 \\
Z-MAC*                & 0.0663 & 0.0325 & 0.0552 & 19.5523 & 18764.5 & 0.5731 & 0.8611 & 0.0092 \\
TRAMA*                & 0.0371 & 0.0182 & 0.0153 & 18.8684 & 18644.7 & 0.6077 & 0.5278 & 0.0175 \\
Bandit-Index*         & 0.0343 & 0.0168 & 0.0075 & 20.3143 & 18662.2 & 0.8422 & 0.6944 & 0.0105 \\
\addlinespace[2pt]
PPO-Base                   & 0.2710 & 0.1329 & 0.0452 & 17.1105 & 16062.1 & 0.2765 & 0.6944 & 0.0829 \\
HC-PPO                     & \textbf{0.2770} & \textbf{0.1356} & \textbf{0.0464} & \textbf{16.6721} & \textbf{15824.8} & 0.4921 & 0.6944 & 0.0847 \\
\textbf{DASD-PPO w/o Mask} & 0.2685 & 0.1319 & 0.0390 & 16.8423 & \textbf{15113.0} & 0.2784 & \textbf{0.5000} & \textbf{0.1343} \\
\textbf{DASD-PPO}          & 0.2720 & 0.1334 & 0.0388 & 18.5290 & 16039.4 & 0.4433 & 0.6944 & 0.0832 \\
\bottomrule
\end{tabular}%
}
\end{table}

\begin{table}[!htbp]
\centering
\caption{UAV 稀疏主场景（\texttt{N12\_grid\_manet\_uav\_sparse}）静态对比，3 seeds 平均}
\label{tab:static-uav-sparse}
\setlength{\tabcolsep}{4pt}
\resizebox{\textwidth}{!}{%
\begin{tabular}{@{}lrrrrrrrr@{}}
\toprule
\textbf{方法} &
  \textbf{Goodput/slot}\(\uparrow\) &
  \textbf{PDR}\(\uparrow\) &
  \textbf{MAC eff.}\(\uparrow\) &
  \textbf{Queue slope}\(\downarrow\) &
  \textbf{\(W_t\) P95}\(\downarrow\) &
  \textbf{Jain 尾 100}\(\uparrow\) &
  \textbf{Starv.}\(\downarrow\) &
  \textbf{Eff. Svc.}\(\uparrow\) \\
\midrule
plain TDMA            & 0.0493 & 0.0242 & 0.0411 & 14.2274 & 18534.9 & 0.6533 & 0.8333 & 0.0082 \\
Heuristic             & 0.2723 & 0.1335 & 0.0227 & 10.2692 & 15539.0 & 0.2803 & 0.5556 & 0.1209 \\
Greedy-STDMA-2hop     & 0.2661 & 0.1305 & 0.0224 & 10.3666 & 15758.1 & 0.3006 & 0.4444 & 0.1479 \\
Traffic-Adaptive TDMA & 0.0852 & 0.0418 & 0.0071 & 12.2306 & 17742.7 & 0.8165 & 0.6667 & 0.0284 \\
Z-MAC*                & 0.0544 & 0.0267 & 0.0454 & 13.4321 & 18451.8 & 0.7542 & 0.8056 & 0.0106 \\
TRAMA*                & 0.0728 & 0.0357 & 0.0100 & 13.9033 & 17903.3 & 0.7415 & 0.5833 & 0.0303 \\
Bandit-Index*         & 0.0567 & 0.0278 & 0.0124 & 13.2741 & 18126.0 & 0.8180 & 0.6944 & 0.0173 \\
\addlinespace[2pt]
PPO-Base                   & 0.2423 & 0.1188 & 0.0403 & 10.2136 & 14884.8 & 0.4050 & 0.3889 & 0.1480 \\
HC-PPO                     & 0.2430 & 0.1191 & 0.0406 & 9.9872 & 14682.9 & 0.3690 & 0.5556 & 0.1079 \\
\textbf{DASD-PPO w/o Mask} & 0.2598 & 0.1276 & 0.0395 & \textbf{9.5979} & \textbf{13528.3} & 0.3304 & 0.2778 & 0.1876 \\
\textbf{DASD-PPO}          & \textbf{0.2683} & \textbf{0.1315} & 0.0313 & 12.0047 & 14886.2 & 0.3555 & \textbf{0.2500} & \textbf{0.2012} \\
\bottomrule
\end{tabular}%
}
\end{table}

\sloppy
两张表呈现的结果可以从三个维度解读。
\emph{首先}，学习型算法整体显著优于
paper-inspired 论文启发式近似：
PPO-Base 与 HC-PPO 在行人场景的 Goodput/slot 较
Z-MAC*、TRAMA* 与 Bandit-Index* 高约 4--8 倍；
UAV 稀疏场景中三个 paper-inspired 方法 Goodput/slot
均不足 0.08，而所有学习型算法均超过 0.24。
这一差距证实了在动态拓扑下，
基于在线策略梯度的调度比固定权值仲裁更有竞争力。
\fussy

\emph{其次}，行人场景中 HC-PPO 在七项指标中的四项
（Goodput/slot、PDR、MAC 效率、\(W_t\) P95、Queue slope）取得最优，
表明启发式偏离惩罚对密集拓扑确实有正向收益；
DASD-PPO 在该场景下的吞吐与 PDR 与 HC-PPO 基本持平
（Goodput/slot 0.2720 vs 0.2770，相对回退 \(\sim 1.8\%\)），
但 Queue slope 由 16.67 升至 18.53，说明
两跳掩膜与服务债务在邻居密度较高的拓扑里增加了队列累积压力。
因此 DASD-PPO 在行人场景\emph{并未}全面优于 HC-PPO，
不应被宣称为该场景下的主方法。

\emph{第三}，UAV 稀疏场景中 DASD-PPO 的特征则相反。
其 Goodput/slot（0.2683）相较 HC-PPO（0.2430）提升约 \(10.4\%\)，
接近 Heuristic（0.2723）；
其 Starvation 比例由 HC-PPO 的 0.5556 与 Heuristic 的 0.5556
显著下降至 \textbf{0.2500}，是全表所有算法中最低。
代价是 Queue slope 由 HC-PPO 的 9.99 升至 12.00。
该现象与算法设计一致：在稀疏拓扑下两跳安全集合较小，
服务债务机制能更精准地为长期未服务节点分配安全候选；
而密度自适应控制让掩膜在密度更高的行人场景中部分退化，
减少了过严约束带来的可行域损失。

\paragraph{两跳掩膜的边际贡献：意外的负向证据。}
将 DASD-PPO w/o Mask 与 DASD-PPO 直接对比，
可以读出两跳掩膜在静态主场景中的真实边际贡献：
\emph{行人场景}下，去掉掩膜反而把 Queue slope 从 \(18.53\) 降到 \(16.84\)（\(-9.1\%\)）、
\(W_t\) P95 从 \(16039\) 降到 \(\mathbf{15113}\)（全表最低）、
Starvation 从 \(0.6944\) 降到 \(\mathbf{0.5000}\)，
吞吐与 PDR 仅微降 \(\sim 1.3\%\)；
\emph{UAV 稀疏场景}下，去掉掩膜把 Queue slope 从 \(12.00\) 降到 \(\mathbf{9.60}\)（\(-20\%\)）、
\(W_t\) P95 从 \(14886\) 降到 \(\mathbf{13528}\)，
Goodput 从 \(0.2683\) 降到 \(0.2598\)（\(-3.2\%\)）、
Starvation 仅由 \(0.2500\) 微升至 \(0.2778\)。
这一对比颠覆了"掩膜是 DASD-PPO 性能核心"的设计直觉——
两跳掩膜在密集拓扑中收缩可行域、抬升队列累积压力，
在稀疏拓扑中则与服务债务推力争夺有限的安全候选，
其唯一明确的净收益是 UAV 场景的 Starvation 微降。
\Cref{subsec:ablation} 与 \Cref{sec:discussion} 将进一步用动态扰动数据验证该结论。

\paragraph{有效服务率：联合解读公平性与吞吐。}
直接对比 Jain 尾 100 容易被低吞吐方法的"伪公平"误导
（Traffic-Adaptive TDMA Jain \(0.890\) 实由 \(0.917\) 的饥饿率支撑），
本文以\emph{有效服务率} \((1-\mathrm{Starv})\times\mathrm{Goodput}\) 作复合排序：
该指标同时奖励"参与的节点多"与"参与节点的吞吐高"，
在所有 \(11\) 种方法上构成无歧义偏序。
行人场景中 \textbf{DASD-PPO w/o Mask 取得全表最高的 \(0.1343\)}，
是 HC-PPO（\(0.0847\)）的 \(1.58\) 倍、是 Bandit-Index*（\(0.0105\)）的 \(12.8\) 倍；
UAV 稀疏场景中 DASD-PPO（\(\mathbf{0.2012}\)）与 DASD-PPO w/o Mask（\(0.1876\)）
分别排名第一与第二，
显著高于 Greedy-STDMA-2hop（\(0.1479\)）与 Heuristic（\(0.1209\)）。
有效服务率视角下，DASD-PPO 与 DASD-PPO w/o Mask 的差距在两场景间互有胜负，
进一步说明两跳掩膜的取舍是场景依赖而非性能必需。

综合两个场景，DASD-PPO 的合理定位是\emph{在稀疏拓扑场景下显著改善公平性与饥饿，
保持吞吐接近 Heuristic 与 HC-PPO 的水平}；
而 DASD-PPO w/o Mask 在行人场景下提供\emph{更低队列压力与更低饥饿}，
在 UAV 稀疏场景下\emph{以微弱吞吐回退换取更短 \(W_t\) P95 与队列稳定}。
两者均不应被宣称为"全场景全指标最优"。

\subsection{动态扰动结果}
\label{subsec:dynamic-results}

\Cref{fig:dyn-pdr-trend}--\Cref{fig:dyn-phase-bars} 给出四个动态扰动场景下
全部九种算法的按帧指标趋势与三阶段统计，
\Cref{tab:dyn-phase} 汇总了三阶段平均 PDR 与 throughput。
新增的 \textbf{DASD-PPO w/o Mask} 行用于在动态扰动维度上单独定位两跳掩膜的边际贡献。

\begin{figure}[!htbp]
\centering
\includegraphics[width=\linewidth]{../results/dynamic_paper_extra_with_nomask_20260524/dyn_pdr_trend.pdf}
\caption{四个动态扰动场景下九种算法的按帧 PDR 趋势。
红色虚线为扰动起点（frame \(=20000\)），
绿色虚线为恢复起点（frame \(=40000\)）；
Node Dropout 场景无恢复阶段。
曲线为同方法 3 个 seed 的均值与 500-帧滑动窗口平滑。}
\label{fig:dyn-pdr-trend}
\end{figure}

\begin{figure}[!htbp]
\centering
\includegraphics[width=\linewidth]{../results/dynamic_paper_extra_with_nomask_20260524/dyn_jain_trend.pdf}
\caption{按帧 Jain 公平指数（发送公平性 \texttt{jain\_tx}）趋势。
学习型算法（实线）在扰动后的公平性恢复速度普遍快于
paper-inspired 方法（点线）。}
\label{fig:dyn-jain-trend}
\end{figure}

\begin{figure}[!htbp]
\centering
\includegraphics[width=\linewidth]{../results/dynamic_paper_extra_with_nomask_20260524/dyn_throughput_trend.pdf}
\caption{按帧每帧 throughput 趋势（单位：分组/帧）。
DASD-PPO 在 Node Rejoin 恢复段的 throughput 取得 2.6072，
高于 HC-PPO（2.4310）与 PPO-Base（2.5613）。}
\label{fig:dyn-throughput-trend}
\end{figure}

\begin{figure}[!htbp]
\centering
\includegraphics[width=\linewidth]{../results/dynamic_paper_extra_with_nomask_20260524/dyn_phase_bars.pdf}
\caption{四个动态扰动场景下九种算法在 Pre/Perturb/Recovery 三阶段
平均 PDR 的分组条形对照。
Node Dropout 场景无 Recovery 阶段，已在图内标注。}
\label{fig:dyn-phase-bars}
\end{figure}

\begin{table}[!htbp]
\centering
\caption{动态扰动主表：四场景 \(\times\) 八方法 \(\times\) 三阶段平均 PDR 与 throughput，
3 seeds 平均。Node Dropout 场景的 Recovery 列按设计为空。}
\label{tab:dyn-phase}
\scriptsize
\setlength{\tabcolsep}{3pt}
\resizebox{\textwidth}{!}{%
\begin{tabular}{@{}llrrrrrr@{}}
\toprule
\textbf{场景} & \textbf{方法} &
  \textbf{PDR pre}\(\uparrow\) & \textbf{PDR perturb}\(\uparrow\) & \textbf{PDR rec.}\(\uparrow\) &
  \textbf{TP pre}\(\uparrow\) & \textbf{TP perturb}\(\uparrow\) & \textbf{TP rec.}\(\uparrow\) \\
\midrule
\multirow{9}{*}{Edge Toggle}
  & DASD-PPO            & 0.1444 & 0.1220 & 0.1312 & 2.7946 & 2.3609 & 2.5386 \\
  & DASD-PPO w/o Mask   & 0.1364 & \textbf{0.1415} & 0.1337 & 2.6357 & \textbf{2.7355} & 2.5794 \\
  & HC-PPO              & 0.1590 & 0.1384 & 0.1390 & 3.0758 & 2.6804 & 2.6841 \\
  & PPO-Base            & 0.1446 & 0.1403 & 0.1346 & 2.7957 & 2.7104 & 2.5953 \\
  & Z-MAC*              & 0.0407 & 0.0326 & 0.0268 & 0.7872 & 0.6298 & 0.5202 \\
  & TRAMA*              & 0.0195 & 0.0204 & 0.0224 & 0.3744 & 0.3960 & 0.4313 \\
  & Bandit-Index*       & 0.0205 & 0.0169 & 0.0174 & 0.3961 & 0.3273 & 0.3369 \\
  & Greedy-STDMA-2hop   & 0.1092 & 0.0704 & 0.0692 & 2.1082 & 1.3603 & 1.3395 \\
  & Heuristic           & 0.0844 & 0.0691 & 0.0695 & 1.6374 & 1.3403 & 1.3435 \\
\midrule
\multirow{9}{*}{Bridge Break}
  & DASD-PPO            & 0.1409 & 0.1357 & 0.1329 & 2.7206 & 2.6236 & 2.5685 \\
  & DASD-PPO w/o Mask   & 0.1479 & 0.1352 & 0.1447 & 2.8608 & 2.6144 & \textbf{2.8005} \\
  & HC-PPO              & 0.1474 & 0.1334 & 0.1353 & 2.8491 & 2.5781 & 2.6143 \\
  & PPO-Base            & 0.1563 & 0.1417 & 0.1345 & 3.0206 & 2.7345 & 2.5974 \\
  & Z-MAC*              & 0.0407 & 0.0326 & 0.0268 & 0.7872 & 0.6298 & 0.5202 \\
  & TRAMA*              & 0.0195 & 0.0217 & 0.0183 & 0.3744 & 0.4205 & 0.3539 \\
  & Bandit-Index*       & 0.0205 & 0.0169 & 0.0175 & 0.3961 & 0.3272 & 0.3378 \\
  & Greedy-STDMA-2hop   & 0.1092 & 0.0704 & 0.0692 & 2.1082 & 1.3602 & 1.3395 \\
  & Heuristic           & 0.0844 & 0.0691 & 0.0695 & 1.6374 & 1.3403 & 1.3435 \\
\midrule
\multirow{9}{*}{Node Rejoin}
  & DASD-PPO            & 0.1461 & 0.0134 & 0.1348 & 2.8256 & 0.1899 & \textbf{2.6072} \\
  & DASD-PPO w/o Mask   & 0.1443 & 0.0135 & 0.1315 & 2.7885 & 0.1916 & 2.5415 \\
  & HC-PPO              & 0.1473 & 0.0153 & 0.1255 & 2.8439 & 0.2188 & 2.4310 \\
  & PPO-Base            & 0.1447 & 0.0121 & 0.1321 & 2.7973 & 0.1711 & 2.5613 \\
  & Z-MAC*              & 0.0407 & 0.0055 & 0.0317 & 0.7872 & 0.0790 & 0.6125 \\
  & TRAMA*              & 0.0195 & 0.0031 & 0.0200 & 0.3744 & 0.0431 & 0.3864 \\
  & Bandit-Index*       & 0.0205 & 0.0025 & 0.0179 & 0.3961 & 0.0343 & 0.3451 \\
  & Greedy-STDMA-2hop   & 0.1092 & 0.0142 & 0.0666 & 2.1082 & 0.2040 & 1.2851 \\
  & Heuristic           & 0.0844 & 0.0136 & 0.0695 & 1.6374 & 0.1929 & 1.3415 \\
\midrule
\multirow{9}{*}{Node Dropout}
  & DASD-PPO            & 0.1454 & 0.0108 & ---    & 2.8093 & 0.1530 & --- \\
  & DASD-PPO w/o Mask   & 0.1491 & 0.0115 & ---    & 2.8876 & 0.1629 & --- \\
  & HC-PPO              & 0.1530 & 0.0112 & ---    & 2.9558 & 0.1595 & --- \\
  & PPO-Base            & 0.1482 & 0.0173 & ---    & 2.8664 & 0.2476 & --- \\
  & Z-MAC*              & 0.0407 & 0.0055 & ---    & 0.7872 & 0.0790 & --- \\
  & TRAMA*              & 0.0195 & 0.0031 & ---    & 0.3744 & 0.0431 & --- \\
  & Bandit-Index*       & 0.0205 & 0.0025 & ---    & 0.3961 & 0.0343 & --- \\
  & Greedy-STDMA-2hop   & 0.1092 & 0.0142 & ---    & 2.1082 & 0.2040 & --- \\
  & Heuristic           & 0.0844 & 0.0136 & ---    & 1.6374 & 0.1929 & --- \\
\bottomrule
\end{tabular}%
}
\end{table}

四个场景的结果不应被概括为"DASD-PPO 全场景领先"。
相反，三种学习型算法在不同扰动模式下各有所长，
共同呈现出"\emph{学习型方法整体显著鲁棒、内部存在场景化分工}"的结构。

在 \textbf{Edge Toggle} 场景中，PDR 衰减最小的是 PPO-Base
（pre \(\rightarrow\) perturb 仅 \(-2.99\%\)），
其次是 HC-PPO（\(-13.0\%\)），
而 DASD-PPO 反而衰减最大（\(-15.5\%\)，pre \(0.1444\rightarrow\) perturb \(0.1220\)）。
进一步诊断表明 DASD-PPO 在 pre 阶段的 sent/candidate 比例（\(\approx 0.586\)）
显著高于 PPO-Base（\(\approx 0.500\)），属于"过申请"现象——
两跳掩膜对该场景中频繁开断的边给出了过于乐观的安全集合，
而服务债务推力放大了这种倾向。
\textbf{DASD-PPO w/o Mask 提供了关键反证}：
关掉掩膜后，pre \(\rightarrow\) perturb 不仅没有衰减，反而 PDR 由 \(0.1364\) 升至 \(\mathbf{0.1415}\)（\(+3.7\%\)），
throughput 由 \(2.636\) 升至 \(\mathbf{2.736}\)（\(+3.8\%\)），
彻底消除了 DASD-PPO 的退化。
这一对照直接证明：\emph{DASD-PPO 在 Edge Toggle 中的衰减完全由两跳掩膜的"快照过期"引起}——
掩膜以 \(t{-}1\) 帧占用为依据屏蔽时隙，而频繁开断的边让该快照在 \(t\) 帧立即失效，
反而把好时隙也屏蔽掉了。

在 \textbf{Bridge Break} 场景中，DASD-PPO 的 PDR 衰减最小（\(-3.72\%\)），
但 pre 阶段 PDR 的绝对值（\(0.1409\)）也最低；
PPO-Base 的 pre 最高（\(0.1563\)）但衰减更大（\(-9.3\%\)）；
\textbf{DASD-PPO w/o Mask 在恢复段取得全表最高 throughput \(\mathbf{2.8005}\)}
（vs DASD-PPO \(2.5685\)，相对提升 \(+9.0\%\); vs HC-PPO \(2.6143\)，提升 \(+7.1\%\)），
说明掩膜的"快照锁定"效应在结构性扰动恢复期同样起反作用——
关掉掩膜后，服务债务能更快地将新拓扑下的安全候选分配给受影响节点。
这一对照进一步弱化了掩膜的必要性：
DASD-PPO 借助服务债务对受影响节点的补偿可以减小退化幅度，
但其代价是 pre 阶段以略低的吞吐换取保守，
且 Recovery 段 throughput 反而被 DASD-PPO w/o Mask 显著超越。

在 \textbf{Node Rejoin} 场景中，四种学习型算法在 perturb 段
PDR 同时跌至 \(0.012\sim 0.015\) 区间——
即新加入节点尚未被任何调度策略所"识别"，
分组几乎全部丢失。
这是当前学习型 MAC 方案的共同局限：
现有策略并未在训练分布中覆盖"新节点上线"的事件，
故无法做出零样本响应。
Recovery 段 DASD-PPO 的 throughput（\(2.6072\)）显著高于
HC-PPO（\(2.4310\)）与 PPO-Base（\(2.5613\)），
DASD-PPO w/o Mask 的 throughput（\(2.5415\)）位居其后，
说明在拓扑稳定后，服务债务是恢复长尾的核心机制，
而两跳掩膜在该场景中提供小幅但可观察的额外贡献（\(+2.5\%\) recovery TP）。
这是本节中两跳掩膜\emph{唯一明确表现为净正向贡献}的扰动场景。

在 \textbf{Node Dropout} 场景中，HC-PPO 的 pre PDR 最高（\(0.1530\)），
\textbf{DASD-PPO w/o Mask 紧随其后取得 \(0.1491\)}（高于 DASD-PPO 的 \(0.1454\) 与 PPO-Base 的 \(0.1482\)）；
所有算法在 perturb 段 PDR 均跌至 \(0.011\sim 0.017\)
（永久脱离节点的分组难以投递）；
由于该场景在仿真器实现中按设计跳过恢复阶段，
所有算法的 Recovery 列在统计上空缺，
不应作为评估依据。

值得单独剖析的是\textbf{Greedy-STDMA-2hop}——它在两个静态主场景里都是传统启发式中的
最强基线（行人 Goodput/slot \(0.1608\)、UAV 稀疏 \(0.2661\)，均居同组之首），
然而在四个动态扰动场景下都呈现典型的"\emph{染色快照过期}"型衰减：
Edge Toggle 与 Bridge Break 中其 PDR 从 pre 阶段的 \(0.1092\) 跌至 perturb 段的 \(0.0704\)
（衰减 \(-35.5\%\)），原因在于两跳贪心染色得到的复用方案本质上是
\emph{当前帧拓扑的一份快照}——一旦边集合切换或关键链路断开，
原方案立即对部分时隙产生新的冲突关系，
而 Greedy-STDMA-2hop 不具备在线再学习能力、
需要等到下一次重新染色才能恢复，期间产生大量碰撞与丢包。
该现象在 \textbf{Node Rejoin} 中更为剧烈：
perturb 段 PDR 跌至 \(0.0142\)，因为新加入节点完全不在原冲突图中、
染色方案对其"不可见"；
Recovery 段仅回到 \(0.0666\)（约为 pre 的 \(61\%\)），
亦说明重新染色后的方案与扰动前相比仍处于次优。
同源对比下，DASD-PPO 在 Node Rejoin 恢复段的 throughput 取得全表最高的 \(2.6072\)，
正是\emph{RL 在线更新机制}相对\emph{离线染色快照}在动态拓扑下结构性优势的直接证据。

paper-inspired 三个基线在所有四个场景的 PDR 全程不超过 \(0.04\)、
throughput 全程不超过 \(0.8\)，
显著低于\emph{任何}学习型方法，
也低于 Greedy-STDMA-2hop 与 Heuristic 两种传统强基线。
这一全方位差距是支撑"在线学习对动态扰动具有结构性鲁棒优势"
最直接的实验证据。
配合 \Cref{fig:dyn-jain-trend} 中 Jain 公平性的恢复速度，
学习型方法即便在 Node Rejoin/Dropout 这种结构性扰动下也仍能
更快重建公平的发送分布。

\subsection{算法演进消融}
\label{subsec:ablation}

为定位 DASD-PPO 中各机制的边际贡献，
本节将四个学习型变体——
\textbf{PPO-Base \(\rightarrow\) HC-PPO \(\rightarrow\) DASD-PPO w/o Mask \(\rightarrow\) DASD-PPO}——
按机制叠加顺序构造四级消融对照
（机制差异见 \Cref{tab:algo-naming} 与 \Cref{fig:algo-evolution}）。
该链条把"两跳掩膜"从原本被打包在 DASD 模块内的隐式贡献剥离出来，
让前三个机制（HC 锚定、服务债务 + 密度自适应 + 请求预算）和掩膜
能被独立测量。

\paragraph{HC 模块（启发式偏离锚定）的贡献。}
PPO-Base \(\rightarrow\) HC-PPO 的差异是 \eqref{eq:heur-deviation} 的软惩罚项。
在行人场景中，HC-PPO 相对 PPO-Base 在 Goodput/slot（\(0.2710\rightarrow 0.2770\)，
\(+2.2\%\)）、PDR（\(0.1329\rightarrow 0.1356\)）、
Queue slope（\(17.11\rightarrow 16.67\)）与 \(W_t\) P95（\(16062\rightarrow 15825\)）四项均小幅改进，
且 Jain 尾 100 从 \(0.28\) 升至 \(0.49\)。
在 UAV 稀疏场景中差距更小（Goodput/slot \(0.2423 \rightarrow 0.2430\)），
说明 HC 模块的主要收益体现在密度较高、启发式可行域较窄的拓扑。

\paragraph{服务债务 + 密度自适应 + 请求预算（无掩膜版）的贡献。}
HC-PPO \(\rightarrow\) DASD-PPO w/o Mask 的差异是 \Cref{subsec:debt}--\Cref{subsec:density}
三节的机制叠加，但 \emph{不含} \Cref{subsec:mask} 的两跳掩膜。
\emph{静态收益}：行人场景 Starvation 由 \(0.6944\) 降至 \(\mathbf{0.5000}\)（\(-28\%\)），
\(W_t\) P95 由 \(15825\) 降至 \(\mathbf{15113}\)（全表最低）；
UAV 稀疏场景 Goodput/slot 由 \(0.2430\) 升至 \(0.2598\)（\(+6.9\%\)），
Queue slope 由 \(9.99\) 降至 \(\mathbf{9.60}\)（全表最低），
Starvation 由 \(0.5556\) 降至 \(0.2778\)（\(-50\%\)），
有效服务率由 \(0.1079\) 升至 \(0.1876\)（\(+74\%\)）。
\emph{动态收益}：Edge Toggle PDR 由 HC-PPO 的 \(-13.0\%\) 衰减改善为 \(+3.7\%\)\emph{反向增益}；
Bridge Break Recovery throughput 由 \(2.6143\) 升至 \(\mathbf{2.8005}\)（\(+7.1\%\)，全表最高）。
该子集已贡献了 DASD-PPO 大部分的公平性与扰动鲁棒性收益。

\paragraph{两跳掩膜（mask）的单独贡献：净增益微弱甚至为负。}
DASD-PPO w/o Mask \(\rightarrow\) DASD-PPO 的差异\emph{仅是} \Cref{subsec:mask} 的两跳掩膜
与配套的 \(\alpha=0.75\) 安全槽偏置。
该步带来：
\emph{正向}——UAV 稀疏场景 Starvation 由 \(0.2778\) 微降至 \(0.2500\)；
Node Rejoin Recovery throughput 由 \(2.5415\) 升至 \(2.6072\)（\(+2.5\%\)）；
有效服务率 UAV 稀疏由 \(0.1876\) 升至 \(0.2012\)（\(+7.2\%\)）。
\emph{负向}——行人场景 Queue slope 由 \(16.84\) 升至 \(18.53\)（\(+10\%\)），
\(W_t\) P95 由 \(15113\) 升至 \(16039\)（\(+6\%\)），
Starvation 由 \(0.5000\) 升至 \(0.6944\)（\(+39\%\)）；
UAV 稀疏 Queue slope 由 \(9.60\) 升至 \(12.00\)（\(+25\%\)），
\(W_t\) P95 由 \(13528\) 升至 \(14886\)（\(+10\%\)）；
\textbf{Edge Toggle PDR perturb 由 \(+3.7\%\) 反向增益重新跌至 \(-15.5\%\) 衰减}；
Bridge Break Recovery throughput 由 \(\mathbf{2.8005}\) 降至 \(2.5685\)（\(-8.3\%\)）。

将所有 \(11\) 种方法 \(\times\,6\) 种场景 \(\times\) 多指标的比较聚合后：
\textbf{两跳掩膜的净贡献仅在 UAV 稀疏与 Node Rejoin 中明确为正，
在行人静态与 Edge Toggle/Bridge Break 中明确为负}。
这与最初将掩膜作为"DASD-PPO 性能核心"的设计直觉相悖：
掩膜的真正定位更接近"防 RL 策略采样到非法动作的硬约束"，
其性能贡献已被服务债务+密度自适应+请求预算这三件套覆盖。
\Cref{sec:discussion} 将基于该消融结果进一步重新定位掩膜的角色。

\section{讨论}
\label{sec:discussion}

\subsection{设计权衡与场景适配}

\Cref{subsec:ablation} 的四级消融揭示了 DASD-PPO 各机制之间的内禀权衡。
\emph{服务债务 + 密度自适应 + 请求预算}
是 DASD-PPO 公平性与扰动鲁棒性收益的真正来源——
该子集（即 DASD-PPO w/o Mask）在两个静态主场景的有效服务率均跻身全表前二，
在 Edge Toggle 与 Bridge Break 上的 PDR 衰减抗性与 Recovery throughput 反优于完整 DASD-PPO。
\emph{两跳掩膜}的角色则需要重新定位：
其原本被设计为"显式拓扑感知归纳偏置"，
实际测量结果表明它在密集拓扑下会与服务债务推力争夺有限可行域，
在高动态边切换场景下因"占用快照"过期而产生反向作用，
仅在 UAV 稀疏与 Node Rejoin 中提供\emph{微弱}的净正贡献
（前者 Starvation \(-0.028\)、后者 Recovery TP \(+2.5\%\)）。
\emph{密度自适应控制}通过 \(\rho_t\) 阈值平滑掩膜启停，
本是对掩膜负面效应的部分缓解机制；
但消融数据表明，简单地\emph{完全关闭}掩膜在多数场景下比"按密度部分开启"更具竞争力。

综合实验，DASD-PPO 系列的合理使用画像分为两条：
\begin{itemize}
\item \textbf{DASD-PPO}（带掩膜）：适合 \emph{UAV 稀疏 + 节点动态加入} 场景——
  其唯一明确赢面是 Starvation 微降与 Node Rejoin Recovery 微强。
\item \textbf{DASD-PPO w/o Mask}：适合 \emph{密集拓扑 + 频繁链路开断} 场景——
  其 Queue slope、\(W_t\) P95、Edge Toggle 衰减抗性、Bridge Break Recovery throughput
  在四个学习型变体中均为最佳或并列最佳。
\end{itemize}
两者均不应被宣称为"全场景全指标最优"。
更进一步地，本文建议将两跳掩膜从"主方法的必备机制"降级为
"\emph{条件触发的可选保护层}"——
仅在网络密度持续低于阈值\emph{且}拓扑近似稳态时启用，
其余场景默认关闭以避免快照过期带来的反向作用。

\subsection{已尝试但未采纳的变体}

\sloppy
工程实验登记中已记录的若干尝试也提供了对设计空间的边界探索：
\textbf{固定 \(W_t\) 阈值的请求闸门}单独使用并不能跨场景成立；
\textbf{显式拉格朗日约束变体}由于约束信号容易饱和到 \(1.0\)，
没有取代 HC-PPO 或 DASD-PPO；
\textbf{多头 Critic / CTDE 变体}
在 UAV 稀疏场景下的 long-gap 改善有限，
但在行人场景下 queue slope 反弹。
这些负向结果在工程实验登记中已留下完整轨迹，
不作为本文主结论，
但提示后续研究在引入更复杂的约束机制时需慎重评估其与已有机制的相互作用。
\fussy

\subsection{局限性}
\label{subsec:limitations}

本文存在以下五点基本局限：
\textbf{(i)} 在 Node Rejoin/Dropout 这种结构性扰动下，
所有学习型算法在 perturb 段都无法做零样本响应，
说明当前 PPO 训练分布不覆盖"新节点上线"事件；
\textbf{(ii)} DASD-PPO 的密度阈值 \(\rho_{\mathrm{sparse}},\rho_{\mathrm{dense}}\)
与掩膜的"条件触发"判据当前仍是手工选定，缺乏在线超参数调节机制；
\textbf{(iii)} 实验仅在 \(N=12\)、\(M=10\) 规模下验证，
更大规模或异构拓扑下的可扩展性尚未在本轮实验中评估；
\textbf{(iv)} 能耗、信道速率敏感性以及 \texttt{sim\_time}\(>20000\)
的长期稳定性\emph{该项未在本轮实验中统计}；
\textbf{(v)} 本文以 DASD-PPO w/o Mask 作为掩膜消融对照变体，
其 \(\alpha=0.5\) 安全槽初始化是与 PPO-Base/HC-PPO 对齐的"纯关掩膜"实现；
若进一步搜索"无掩膜下重新调优的 \(\alpha\) 偏置"，
其相对 DASD-PPO 的优势可能进一步扩大，本文未做该方向的额外消融。

\section{结论与展望}
\label{sec:conclusion}

本文针对动态多跳无线网络的 TDMA 调度问题，
提出了密度感知服务债务 PPO（DASD-PPO）调度框架，
通过\emph{乘数模式策略合成 \(+\) 启发式偏离锚定 \(+\) 两跳安全掩膜
\(+\) 服务债务 \(+\) 请求预算 \(+\) 密度自适应控制}
六项机制将 PPO 学习信号约束在拓扑感知可行域内，
并对长尾节点提供确定性补偿。
配合 LSTM Actor--Critic 主干（每节点独立、约 13.8K 参数）
与逐时隙信用分配，
DASD-PPO 能以在线方式同时学习吞吐与公平性的权衡。

在 OMNeT++ 6.3 平台上的对比实验覆盖了两个静态主场景与四个动态扰动场景，
共比较了四种学习型算法（PPO-Base、HC-PPO、DASD-PPO w/o Mask、DASD-PPO）与
七种启发式/论文启发式近似算法。
其中 DASD-PPO w/o Mask 作为新增的消融对照变体，
关闭两跳掩膜但保留服务债务、请求预算与密度自适应，
其作用是单独定位掩膜的边际贡献。
主要发现可概括为四点：
（i）\emph{学习型方法对动态拓扑具有结构性鲁棒优势}——
PDR、throughput 与 Jain 公平性恢复速度均显著优于 paper-inspired 基线，
有效服务率（\((1-\mathrm{Starv})\times\mathrm{Goodput}\)）上至少高出一个数量级；
（ii）\emph{DASD-PPO 在 UAV 稀疏拓扑下显著改善公平性}——
Starvation 由 HC-PPO 的 \(0.5556\) 降至 \(\mathbf{0.2500}\) 全表最低，
有效服务率 \(\mathbf{0.2012}\) 全表最高；
（iii）\emph{DASD-PPO 的扰动恢复优势可被验证}——
Bridge Break 取得最小 PDR 衰减，
Node Rejoin 恢复段取得最高 throughput（\(2.6072\)）；
（iv）\emph{两跳掩膜并非 DASD-PPO 的性能核心}——
新增的 DASD-PPO w/o Mask 在行人静态、Edge Toggle、Bridge Break 三类场景下
反优于带掩膜的 DASD-PPO（Queue slope/\(W_t\) P95/Edge Toggle 衰减/Bridge Break Recovery TP），
说明掩膜的真正定位是"防 RL 输出非法动作的硬约束"而非性能驱动因子，
其性能贡献已被服务债务+密度自适应+请求预算这三件套覆盖。
实验结果不支持"单一 DASD-PPO 全场景全指标最优"的简单结论；
DASD-PPO 适合 \emph{UAV 稀疏 + 节点动态加入} 场景，
DASD-PPO w/o Mask 适合 \emph{密集拓扑 + 频繁链路开断} 场景。

未来工作将围绕五个方向展开：
（1）将训练分布扩展至节点动态加入/退出事件，
    缓解 Node Rejoin 中策略对结构性扰动的零样本失败；
（2）将两跳掩膜从主方法的必备机制降级为\emph{条件触发的可选保护层}，
    探索基于网络密度与拓扑变更速率的在线启用判据，
    替代当前手工选定的 \(\rho_{\mathrm{sparse}},\rho_{\mathrm{dense}}\) 阈值；
（3）研究 Actor--Critic 解耦与多头 Critic 在更大规模拓扑下的可扩展性；
（4）补全能耗、信道速率敏感性与长仿真时间稳定性等本轮未覆盖的评估维度；
（5）在无掩膜版本下重新调优安全槽偏置 \(\alpha\)，
    探索"调优后无掩膜"相对"完整 DASD-PPO"的进一步性能空间。


% \section*{致谢}

感谢 OMNeT++ 开源社区提供的仿真平台支持。
   % 致谢（如需启用，去掉本行注释）

\section{References}
\printbibliography[heading=none]

\end{document}
